\subsection{Гарантии обычного правила Армихо}\label{sec:armijo}
Рассмотрим известное правило достаточного убывания: в каждой точке
начинаем с $\alpha=1$ и умножаем шаг на фиксированное
$\beta\in(0,1)$ до выполнения
\begin{equation}\label{eq:armijo}
 f(x+\alpha d)\le f(x)-\alpha G(x)/2.
\end{equation}
Используем тот же сертификат \eqref{eq:stop} и сохраняем принятое
значение функции. Это обычный поиск Армихо, не новый алгоритм FW;
для него не нужны $L_0,L_1,D$.

\begin{proposition}[Сертифицированная сложность метода сравнения]
\label{prop:armijo}
В условиях теоремы~\ref{thm:gap_stopping} обозначим через
$\alpha_{\log}(x)$ шаг \eqref{eq:log_step} вместе с предельными случаями.
Поиск конечен, а принятый шаг удовлетворяет
$\alpha\ge\beta\alpha_{\log}(x)$.
При $G(x_0)>\gamma\varepsilon$ число обновлений ограничено:
\[
 N\le\left\lceil\frac1\beta\left[
 \frac{\max\{4,4L_1D^2/\log2\}}\gamma J
 +\frac{16L_0D^2}{(\log2)\gamma^2\varepsilon}
 \right]\right\rceil,
 \qquad J=\left\lceil\log_2\frac{G(x_0)}{\gamma\varepsilon}\right\rceil.
\]
Вне аффинного случая положим
\[
 a_\varepsilon=\min\left\{1,
 \frac{(\log2)\gamma\varepsilon}
 {D^2(L_0+L_1\gamma\varepsilon)}\right\},\qquad
 T_\varepsilon=1+\left\lceil\frac{\log(1/a_\varepsilon)}{\log(1/\beta)}\right\rceil.
\]
Число вычислений градиента и линейного минимума равно $N+1$,
а число вычислений функции не превосходит $1+NT_\varepsilon$.
В аффинном случае для единственного возможного обновления
достаточно одной пробы.
\end{proposition}
\begin{proof}
Доказательство леммы~\ref{lem:descent} применимо к каждому
$0<\alpha\le\alpha_{\log}(x)$: производная модели остаётся
неположительной, а та же оценка интеграла даёт \eqref{eq:armijo}.
Поэтому первая точка геометрической сетки, не превосходящая
$\alpha_{\log}(x)$, будет принята, если раньше не принят больший шаг.
Это доказывает конечность поиска и нижнюю границу шага.
Фазовое доказательство теоремы~\ref{thm:gap_stopping} применяется
с уменьшением гарантированного спуска в $\beta$ раз.
До остановки $G(x)>\gamma\varepsilon$. Монотонность
$G/(L_0+L_1G)$ и \eqref{eq:step_lower} дают
$\alpha_{\log}(x)\ge a_\varepsilon$, откуда достаточно
$T_\varepsilon$ проб. Принятое значение функции используется
повторно при переходе к новой точке.
\end{proof}
При $L_0=0$ сертификат за обновление сокращается не хуже, чем
в $1-\beta\gamma\underline\alpha_0/2$ раза. В условиях
теоремы~\ref{thm:strong_set_linear} аналогичный множитель равен
$1-\beta\gamma\underline\alpha_{\rm sc}/2$.
Следовательно, благоприятная оценка числа обновлений определяется
условием кривизны и не является исключительным свойством
логарифмической формулы. Её особенность --- гарантированный шаг
без отклонённых пробных вычислений функции.
